
\documentclass{article}
\usepackage{colortbl}
\usepackage{makecell}
\usepackage{multirow}
\usepackage{supertabular}

\begin{document}

\newcounter{utterance}

\centering \large Interaction Transcript for game `wordle', experiment `medium\_frequency\_words\_no\_clue\_no\_critic', episode 7 with Qwen3.6{-}35B{-}A3B{-}FP8{-}without{-}reasoning.
\vspace{24pt}

{ \footnotesize  \setcounter{utterance}{1}
\setlength{\tabcolsep}{0pt}
\begin{supertabular}{c@{$\;$}|p{.15\linewidth}@{}p{.15\linewidth}p{.15\linewidth}p{.15\linewidth}p{.15\linewidth}p{.15\linewidth}}
    \# & \multicolumn{2}{c}{Player} && \multicolumn{2}{c}{Game Master} \\
    \hline

    \theutterance \stepcounter{utterance}  
    & & & \multicolumn{4}{p{0.6\linewidth}}{
        \cellcolor[rgb]{0.9,0.9,0.9}{
            \makecell[{{p{\linewidth}}}]{
                \texttt{\tiny{[P1$\langle$GM]}}
                \texttt{Verilen kuralları kullanarak kelime tahmin etmeyi seven bir dil sihirbazısın.} \\
\\ 
\texttt{Wordle’a hoş geldin! Hedef kelimeyi tahmin etmek için altı denemen var: beş küçük harften oluşan, geçerli bir kelime. Her tahmin için kısa bir açıklama sunmak üzere lütfen "açıklama:" ve "tahmin:" etiketlerini kullan.} \\
\\ 
\texttt{Örneğin, tahminin "yemek" ise, yanıtın şöyle olmalı:} \\
\texttt{açıklama: bu yaygın bir beş harfli kelime ve tahminime bu kelimeyle başlıyorum.} \\
\texttt{tahmin: yemek} \\
\\ 
\texttt{Her tahminden sonra yanıtın doğrulanacak ve hangi harflerin doğru olduğunu (yeşil), hangi harflerin doğru ama yanlış konumda olduğunu (sarı) ve hangi harflerin yanlış olduğunu (kırmızı) belirten geri bildirim alacaksın. Bu geri bildirim, bir sonraki tahmininde hangi harfleri dahil edeceğini veya hariç tutacağını belirlemede faydalı olabilir.} \\
\\ 
\texttt{Örneğin, "yemek" için geri bildirim şöyle olabilir:} \\
\texttt{tahmin\_geri\_bildirimi: y<sarı> e<sarı> m<yeşil> e<sarı> k<kırmızı>} \\
\\ 
\texttt{Açıklama, yeni bir tahmine ulaşmak için tahmin\_geri\_bildirimi’in nasıl kullanıldığına dair ayrıntılar içermelidir.} \\
\\ 
\texttt{Başka hiçbir metin üretme.} \\
\\ 
\texttt{İlk tahmininle başlayalım.} \\
            }
        }
    }
     \\ \\

    \theutterance \stepcounter{utterance}  
    & \multicolumn{4}{p{0.6\linewidth}}{
        \cellcolor[rgb]{0.561,0.737,0.561}{
            \makecell[{{p{\linewidth}}}]{
                \texttt{\tiny{[P1$\rangle$GM]}}
                \texttt{açıklama: Kelimeyi bulmak için yaygın kullanılan harfleri içeren, anlamlı bir beş harfli kelimeyle başlıyorum.} \\
\texttt{tahmin: araba} \\
            }
        }
    }
    & & \\ \\

    \theutterance \stepcounter{utterance}  
    & & & \multicolumn{2}{p{0.3\linewidth}}{
        \cellcolor[rgb]{0.9,0.9,0.9}{
            \makecell[{{p{\linewidth}}}]{
                \texttt{\tiny{[GM$|$GM]}}
                \texttt{attempts = 1} \\
\texttt{target\_word = çiçek} \\
\texttt{guess = araba} \\
\texttt{guess\_feedback = a<red> r<red> a<red> b<red> a<red>} \\
            }
        }
    }
    & & \\ \\

    \theutterance \stepcounter{utterance}  
    & & & \multicolumn{4}{p{0.6\linewidth}}{
        \cellcolor[rgb]{0.9,0.9,0.9}{
            \makecell[{{p{\linewidth}}}]{
                \texttt{\tiny{[P1$\langle$GM]}}
                \texttt{tahmin\_geri\_bildirimi: a<kırmızı> r<kırmızı> a<kırmızı> b<kırmızı> a<kırmızı>} \\
\\ 
\texttt{Yanıtını yalnızca bu biçimde ver.} \\
\texttt{açıklama: ayrıntılar} \\
\texttt{tahmin: kelime} \\
            }
        }
    }
     \\ \\

    \theutterance \stepcounter{utterance}  
    & \multicolumn{4}{p{0.6\linewidth}}{
        \cellcolor[rgb]{0.561,0.737,0.561}{
            \makecell[{{p{\linewidth}}}]{
                \texttt{\tiny{[P1$\rangle$GM]}}
                \texttt{açıklama: "araba" kelimesindeki tüm harfler (a, r, b) yanlış çıktı. Bu nedenle bu harfleri tamamen hariç tutuyorum. Kalan harfler arasında en yaygın olanları (e, l, i, n, t, s, o, u) kullanarak yeni bir kombinasyon deneyeceğim.} \\
\texttt{tahmin: liste} \\
            }
        }
    }
    & & \\ \\

    \theutterance \stepcounter{utterance}  
    & & & \multicolumn{2}{p{0.3\linewidth}}{
        \cellcolor[rgb]{0.9,0.9,0.9}{
            \makecell[{{p{\linewidth}}}]{
                \texttt{\tiny{[GM$|$GM]}}
                \texttt{attempts = 2} \\
\texttt{target\_word = çiçek} \\
\texttt{guess = liste} \\
\texttt{guess\_feedback = l<red> i<green> s<red> t<red> e<yellow>} \\
            }
        }
    }
    & & \\ \\

    \theutterance \stepcounter{utterance}  
    & & & \multicolumn{4}{p{0.6\linewidth}}{
        \cellcolor[rgb]{0.9,0.9,0.9}{
            \makecell[{{p{\linewidth}}}]{
                \texttt{\tiny{[P1$\langle$GM]}}
                \texttt{tahmin\_geri\_bildirimi: l<kırmızı> i<yeşil> s<kırmızı> t<kırmızı> e<sarı>} \\
\\ 
\texttt{Yanıtını yalnızca bu biçimde ver.} \\
\texttt{açıklama: ayrıntılar} \\
\texttt{tahmin: kelime} \\
            }
        }
    }
     \\ \\

    \theutterance \stepcounter{utterance}  
    & \multicolumn{4}{p{0.6\linewidth}}{
        \cellcolor[rgb]{0.561,0.737,0.561}{
            \makecell[{{p{\linewidth}}}]{
                \texttt{\tiny{[P1$\rangle$GM]}}
                \texttt{açıklama: "liste" tahmininde 'i' harfi doğru konumda (2. harf). 'e' harfi kelimede var ama yanlış konumda (sarı). 'l', 's', 't' harfleri yanlış. 'e' harfi 5. harf olamaz, 1, 3 veya 4. harf olabilir. 'i' 2. harf olarak sabit. Kalan harfler arasında 'a', 'o', 'u', 'n', 'r', 'd', 'k', 'm', 'p', 'y', 'h', 'g', 'v', 'ç', 'ş', 'c' gibi harfleri deneyebilirim. 'e' harfini 1. veya 3. veya 4. sıraya koyup, 'i'yi 2. sırada tutan bir kelime deneyeceğim. Örneğin 'e'yi 1. sıraya, 'i'yi 2. sıraya alıp, kalan harflerle 'e'yi 4. veya 5. sıraya koymayı deneyebilirim. 'e' 5. harf olamayacağına göre, 1, 3, 4 olabilir. 'e'yi 1. harf yapalım. 'i' 2. harf. 3, 4, 5 için kalan harfler. 'e' 4. harf de olabilir. 'e' 3. harf de olabilir. 'e' 1. harf olursa, 2. 'i'. 3,4,5 için 'a','o','u','n','r','d','k','m','p','y','h','g','v','ç','ş','c' denenecek. 'e'yi 4. harf yapalım. 1,3,5 için harfler. 'e'yi 3. harf yapalım. 1,4,5 için harfler. 'e' 1. harf, 'i' 2. harf, 'e' 4. harf olamaz çünkü 'e' sadece bir kez mi var? Wordle'da harf tekrarları olabilir. 'e' 1. harf, 'i'} \\
            }
        }
    }
    & & \\ \\

    \theutterance \stepcounter{utterance}  
    & & & \multicolumn{2}{p{0.3\linewidth}}{
        \cellcolor[rgb]{0.9,0.9,0.9}{
            \makecell[{{p{\linewidth}}}]{
                \texttt{\tiny{[GM$|$GM]}}
                \texttt{Error: The response should contain the 'tahmin:' keyword exactly once.} \\
            }
        }
    }
    & & \\ \\

    \theutterance \stepcounter{utterance}  
    & & & \multicolumn{2}{p{0.3\linewidth}}{
        \cellcolor[rgb]{0.9,0.9,0.9}{
            \makecell[{{p{\linewidth}}}]{
                \texttt{\tiny{[GM$|$GM]}}
                \texttt{game\_result = ABORT} \\
            }
        }
    }
    & & \\ \\

\end{supertabular}
}

\end{document}
